\documentclass[11pt,letterpaper]{article}
\usepackage{iftex}
\ifXeTeX
  \usepackage{fontspec}   % Tectonic / XeLaTeX
\else
  \usepackage[utf8]{inputenc}   % pdfLaTeX (Overleaf)
  \usepackage[T1]{fontenc}
\fi
\usepackage[spanish,es-noshorthands,es-tabla]{babel}
\usepackage{lmodern}
\usepackage[margin=2.2cm]{geometry}
\usepackage{amsmath}
\usepackage{booktabs}
\usepackage{xcolor}
\usepackage{enumitem}
\usepackage{tcolorbox}
\usepackage{listings}
\usepackage[hidelinks]{hyperref}
% GENERADO por ref/analisis.py: numeros que se citan en el texto. No editar a mano.
\newcommand{\GeoPrac}{0,961}
\newcommand{\GeoPracCyM}{0,959}
\newcommand{\GanaPrac}{5}
\newcommand{\NPrac}{10}
\newcommand{\PeorPrac}{6.837}
\newcommand{\GeoPracInv}{1,042}
\newcommand{\GeoPracInvCyM}{1,042}
\newcommand{\GanaPracInv}{4}
\newcommand{\NPracInv}{10}
\newcommand{\PeorPracInv}{7.535}
\newcommand{\GeoAli}{0,971}
\newcommand{\GeoAliCyM}{0,972}
\newcommand{\GanaAli}{15}
\newcommand{\NAli}{70}
\newcommand{\PeorAli}{5.924}
\newcommand{\GeoAliInv}{1,030}
\newcommand{\GeoAliInvCyM}{1,029}
\newcommand{\GanaAliInv}{16}
\newcommand{\NAliInv}{70}
\newcommand{\PeorAliInv}{6.782}
\newcommand{\WPrac}{6.636}
\newcommand{\CyMPrac}{6.497}
\newcommand{\OrdenesPrac}{2}
\newcommand{\ExpoPrac}{100}
\newcommand{\WPracInv}{15.280}
\newcommand{\CyMPracInv}{15.348}
\newcommand{\OrdenesPracInv}{2}
\newcommand{\ExpoPracInv}{100}
\newcommand{\TbW}{9.995}
\newcommand{\TbOrdenes}{1}
\newcommand{\TbPerdida}{5}
\newcommand{\TbCyM}{6.695}
\newcommand{\TbPeor}{12}
\newcommand{\TbAprobados}{4/4}
\newcommand{\TecCincoRazon}{1,002}
\newcommand{\TecCincoExpo}{100}
\newcommand{\TecCincoGeo}{0,932}
\newcommand{\TecHoraRazon}{1,028}
\newcommand{\TecHoraExpo}{99}
\newcommand{\TecHoraGeo}{1,540}
\newcommand{\TecRondasRazon}{1,004}
\newcommand{\TecRondasExpo}{100}
\newcommand{\TecRondasGeo}{1,085}
\newcommand{\TecSubeRazon}{1,007}
\newcommand{\TecSubeExpo}{100}
\newcommand{\TecSubeGeo}{1,462}
\newcommand{\TecBajaRazon}{1,002}
\newcommand{\TecBajaExpo}{99}
\newcommand{\TecBajaGeo}{0,748}
\newcommand{\PracMin}{72}
\newcommand{\PracMax}{190}
\newcommand{\PracIni}{151}
\newcommand{\PracFin}{98}
\newcommand{\PracVR}{0,53}
\newcommand{\PracAC}{-0,11}
\newcommand{\CiclosPromSw}{406,6}
\newcommand{\CiclosPeorSw}{443}
\newcommand{\ICPromSw}{180,1}
\newcommand{\CPISw}{2,26}
\newcommand{\PalabrasSw}{159}
\newcommand{\TickCeroSw}{1.458}
\newcommand{\MopsSw}{1,32}
\newcommand{\MIPSSw}{11,07}
\newcommand{\PctRtresSw}{19,6}
\newcommand{\SpeedupSw}{1,00}
\newcommand{\CiclosPromShift}{382,6}
\newcommand{\CiclosPeorShift}{419}
\newcommand{\ICPromShift}{168,1}
\newcommand{\CPIShift}{2,28}
\newcommand{\PalabrasShift}{118}
\newcommand{\TickCeroShift}{1.376}
\newcommand{\MopsShift}{1,40}
\newcommand{\MIPSShift}{10,98}
\newcommand{\PctRtresShift}{18,5}
\newcommand{\SpeedupShift}{1,06}
\newcommand{\CiclosPromMul}{133,4}
\newcommand{\CiclosPeorMul}{139}
\newcommand{\ICPromMul}{57,6}
\newcommand{\CPIMul}{2,32}
\newcommand{\PalabrasMul}{121}
\newcommand{\TickCeroMul}{1.213}
\newcommand{\MopsMul}{4,02}
\newcommand{\MIPSMul}{10,80}
\newcommand{\PctRtresMul}{6,1}
\newcommand{\SpeedupMul}{3,05}
\newcommand{\CiclosPromMulShift}{109,4}
\newcommand{\CiclosPeorMulShift}{115}
\newcommand{\ICPromMulShift}{45,6}
\newcommand{\CPIMulShift}{2,40}
\newcommand{\PalabrasMulShift}{80}
\newcommand{\TickCeroMulShift}{1.131}
\newcommand{\MopsMulShift}{4,91}
\newcommand{\MIPSMulShift}{10,42}
\newcommand{\PctRtresMulShift}{5,1}
\newcommand{\SpeedupMulShift}{3,72}
\newcommand{\NTicks}{320.000}
\newcommand{\OpsTick}{21,5}
\newcommand{\BytesTick}{12}
\newcommand{\IntensidadOp}{1,79}
\newcommand{\Presupuesto}{2.262}
\newcommand{\MulPorTick}{4,5}
\newcommand{\FracMul}{72,0}
\newcommand{\FracDup}{8,4}
\newcommand{\CiclosMulSw}{65}
\newcommand{\AmdahlFShift}{8,4}
\newcommand{\AmdahlSShift}{3,4}
\newcommand{\AmdahlPredShift}{1,06}
\newcommand{\AmdahlMedShift}{1,06}
\newcommand{\AmdahlMaxShift}{1,09}
\newcommand{\AmdahlFMul}{72,0}
\newcommand{\AmdahlSMul}{6,5}
\newcommand{\AmdahlPredMul}{2,56}
\newcommand{\AmdahlMedMul}{3,05}
\newcommand{\AmdahlMaxMul}{3,57}
\newcommand{\AmdahlFMulShift}{80,3}
\newcommand{\AmdahlSMulShift}{6,0}
\newcommand{\AmdahlPredMulShift}{3,02}
\newcommand{\AmdahlMedMulShift}{3,72}
\newcommand{\AmdahlMaxMulShift}{5,09}
\newcommand{\LCOriginal}{1191}
\newcommand{\LCPctOriginal}{93,0}
\newcommand{\FmaxOriginal}{42,5--45,5}
\newcommand{\LCSw}{985}
\newcommand{\LCPctSw}{77,0}
\newcommand{\FmaxSw}{45,7--46,6}
\newcommand{\LCShift}{1253}
\newcommand{\LCPctShift}{97,9}
\newcommand{\FmaxShift}{41,3--43,3}
\newcommand{\LCMul}{1185}
\newcommand{\LCPctMul}{92,6}
\newcommand{\FmaxMul}{46,7--47,4}
\newcommand{\LCMulShift}{1444}
\newcommand{\LCPctMulShift}{112,8}
\newcommand{\AhorroRAM}{206}
\newcommand{\LCLibres}{295}
\newcommand{\CostoMul}{200}
\newcommand{\CostoShift}{268}
\newcommand{\ExtGeo}{2,005}
\newcommand{\ExtMediana}{2,391}
\newcommand{\ExtPcinco}{1,068}
\newcommand{\ExtPeor}{0,970}
\newcommand{\ExtPierde}{0,7}
\newcommand{\ExtGanaCym}{100}
\newcommand{\ExtDDmed}{3.465}
\newcommand{\ExtDDpnoventa}{9.117}
\newcommand{\ExtExpo}{51}
\newcommand{\ExtCortaGeo}{1,070}
\newcommand{\ExtCortaMediana}{1,000}
\newcommand{\ExtCortaPcinco}{0,979}
\newcommand{\ExtCortaPeor}{0,899}
\newcommand{\ExtCortaPierde}{9,2}
\newcommand{\ExtCortaGanaCym}{65}
\newcommand{\ExtCortaDDmed}{0}
\newcommand{\ExtCortaDDpnoventa}{2.144}
\newcommand{\ExtCortaExpo}{26}
\newcommand{\TendGeo}{1,111}
\newcommand{\TendMediana}{1,170}
\newcommand{\TendPcinco}{0,476}
\newcommand{\TendPeor}{0,400}
\newcommand{\TendPierde}{41,0}
\newcommand{\TendGanaCym}{34}
\newcommand{\TendDDmed}{6.404}
\newcommand{\TendDDpnoventa}{11.156}
\newcommand{\TendExpo}{100}
\newcommand{\TendCortaGeo}{1,015}
\newcommand{\TendCortaMediana}{1,012}
\newcommand{\TendCortaPcinco}{0,747}
\newcommand{\TendCortaPeor}{0,468}
\newcommand{\TendCortaPierde}{47,7}
\newcommand{\TendCortaGanaCym}{25}
\newcommand{\TendCortaDDmed}{1.862}
\newcommand{\TendCortaDDpnoventa}{4.233}
\newcommand{\TendCortaExpo}{100}
\newcommand{\CymGeo}{1,115}
\newcommand{\CymMediana}{1,158}
\newcommand{\CymPcinco}{0,478}
\newcommand{\CymPeor}{0,399}
\newcommand{\CymPierde}{42,2}
\newcommand{\CymCortaGeo}{1,016}
\newcommand{\CymCortaMediana}{1,012}
\newcommand{\CymCortaPcinco}{0,744}
\newcommand{\CymCortaPeor}{0,461}
\newcommand{\CymCortaPierde}{47,7}
\newcommand{\NVentanas}{2.732}
\newcommand{\ExtGeoDosmilveintidos}{1,784}
\newcommand{\ExtPierdeDosmilveintidos}{1,4}
\newcommand{\TendPierdeDosmilveintidos}{61}
\newcommand{\ExtGeoDosmilveintitres}{2,139}
\newcommand{\ExtPierdeDosmilveintitres}{0,4}
\newcommand{\TendPierdeDosmilveintitres}{30}
\newcommand{\ExtGeoDosmilveinticuatro}{2,143}
\newcommand{\ExtPierdeDosmilveinticuatro}{0,0}
\newcommand{\TendPierdeDosmilveinticuatro}{30}
\newcommand{\MCN}{1.000}
\newcommand{\MCPeor}{0,861}
\newcommand{\MetaW}{23.215}
\newcommand{\MetaWTend}{11.060}
\newcommand{\MetaCym}{9.695}
\newcommand{\MetaWCorta}{9.920}
\newcommand{\MetaWTendCorta}{11.375}
\newcommand{\MetaCymCorta}{11.375}
\newcommand{\MetaTickMax}{10.965}
\newcommand{\MetaTickMin}{17.084}
\newcommand{\ExtPalabras}{63}
\newcommand{\ExtCiclosProm}{12,2}
\newcommand{\ExtCiclosPeor}{22}
\newcommand{\ExtIC}{5,5}
\newcommand{\ExtICPeor}{10}
\newcommand{\ExtCPI}{2,24}
\newcommand{\ExtMops}{5,19}
\newcommand{\ExtMIPS}{11,16}
\newcommand{\ExtOpsTick}{2,5}
\newcommand{\ExtBytesTick}{8}
\newcommand{\ExtNTicks}{1.039.948}
\newcommand{\ExtPctRtres}{1,0}
\newcommand{\ExtSpeedupMul}{11}
\newcommand{\ExtSpeedupSw}{33}
\newcommand{\TbExtW}{9.995}
\newcommand{\TbExtOrdenes}{1}
\newcommand{\TbExtPerdida}{5}
\newcommand{\TbExtCyM}{6.695}
\newcommand{\TbExtPeor}{12}
\newcommand{\TbExtAprobados}{4/4}
\newcommand{\TbMetaW}{23.215}
\newcommand{\TbMetaOrdenes}{8}
\newcommand{\TbMetaPerdida}{3.296}
\newcommand{\TbMetaCyM}{9.695}
\newcommand{\TbMetaPeor}{12}
\newcommand{\TbMetaAprobados}{4/4}
\newcommand{\RefW}{23.215}
\newcommand{\RefOrdenes}{8}
\newcommand{\RefComisiones}{40}
\newcommand{\RefPerdida}{3.296}
\newcommand{\RefCymPerdida}{9.300}
\newcommand{\ExtTiempoProm}{0,49}
\newcommand{\ExtTiempoPeor}{0,88}
\newcommand{\ExtTickCero}{16}
\newcommand{\AEntra}{3}
\newcommand{\ASale}{1}
\newcommand{\CSube}{0}
\newcommand{\CBaja}{24}
\newcommand{\BExceso}{12}


\newtcolorbox{checkpoint}{colback=green!5,colframe=green!45!black,title=Checkpoint,fonttitle=\bfseries\small}
\lstset{basicstyle=\ttfamily\footnotesize,columns=fullflexible,frame=single,framerule=0.3pt,
  rulecolor=\color{gray!50}}
\setlist{itemsep=2pt}
\newcounter{preg}
\newcommand{\preg}[1]{\stepcounter{preg}\par\medskip\noindent\textbf{P\thepreg. #1}\par\noindent}

\title{\textbf{Pochoco Trader: guía para la defensa}\\[2pt]\large Proyecto 4, Arquitectura de Computadores}
\author{Para Martín, Josefa y Felipe}
\date{17 de noviembre de 2026 (algoritmo final)}

\begin{document}
\maketitle

\noindent Esta guía tiene tres partes. La \textbf{A} arma el proyecto desde cero, pieza por pieza, con
un checkpoint al final de cada paso: si no puedes responder el checkpoint, vuelve a leer el paso.
La \textbf{B} son las preguntas que probablemente hagan en la defensa, con su respuesta. La
\textbf{C} son los números que conviene saberse. Todo lo que se cita está en el repositorio.

\tableofcontents

%=============================================================================
\section*{Parte A. El proyecto desde cero}
\addcontentsline{toc}{section}{Parte A. El proyecto desde cero}

\subsection*{Paso 1. El juego}
\addcontentsline{toc}{subsection}{Paso 1. El juego}
\begin{enumerate}
\item La maestra baja CS\_n, manda el precio $p_n$ (1 a 255) y \emph{en la misma transacción} lee
  por MISO el byte que estaba en DECISION.
\item Ese byte lo escribimos después de recibir $p_{n-1}$. La maestra lo ejecuta al precio $p_n$.
\item Byte: bits 7:6 acción (\texttt{01} comprar, \texttt{10} vender), bits 5:0 cantidad $q \le 63$.
\item Si no alcanzamos a escribir antes del siguiente tick, la maestra \textbf{repite} el byte que estaba.
\item La maestra \textbf{recorta} la orden: compra lo que alcance con la caja, vende lo que haya. Una
  orden que queda en 0 no se ejecuta y no paga los \$5 de comisión (tarea \texttt{ejecutar} del testbench).
\item Ronda 3: $T = 0{,}1$\,ms = 2500 ciclos. Menos los 9,5\,$\mu$s del SPI quedan \Presupuesto{} ciclos.
\end{enumerate}
\begin{checkpoint}
¿Qué pasa si mandamos ``comprar 63'' y la caja alcanza para 13 acciones? (Compra 13 y cobra una comisión.
Si no alcanza para ninguna, no pasa nada y no cobra.)
\end{checkpoint}

\subsection*{Paso 2. Cómo se arman las series de la competencia}
\addcontentsline{toc}{subsection}{Paso 2. Cómo se arman las series de la competencia}
\begin{enumerate}
\item El 17-nov llegó \texttt{tb/practica\_meta.csv}: 20.000 minutos de Meta (feb--may 2022),
  normalizados entre 100 y 255. En la competencia: períodos ocultos de Meta, Apple, Netflix, AMD o Tesla
  (2022--2024), normalizados igual.
\item Se bajaron los precios reales de 1 minuto de las 5 empresas (Alpaca, vía Hugging Face).
\item La ventana de Meta que parte el 23-feb-2022 a las 15:20 (horario regular), llevada a $[100,255]$
  y redondeada, da la práctica \textbf{exacta}: 20.000 de 20.000 precios iguales.
\item Fórmula: $p = \operatorname{round}(100 + 155\,(x - \min x)/(\max x - \min x))$, con el mínimo y
  el máximo de \textbf{toda} la ventana.
\item Así se pueden armar todas las series posibles: \NVentanas{} ventanas para evaluar
  (\texttt{ref/minuto.py}, \texttt{ref/competencia.py}).
\end{enumerate}
\begin{checkpoint}
¿Por qué el mínimo y el máximo de toda la ventana le dan información al agente? (Porque incluyen el
futuro: si hoy el precio vale 100, ya se sabe que nunca va a bajar de ahí.)
\end{checkpoint}

\subsection*{Paso 3. El algoritmo final (``extremos'')}
\addcontentsline{toc}{subsection}{Paso 3. El algoritmo final}
\begin{enumerate}
\item Dos hechos de toda ventana: el precio queda entre 100 y 255, y el 100 y el 255 aparecen sí o sí.
\item \textbf{Regla 1}: si $p \le$ piso, dentro (\texttt{0x7F}, comprar 63).
\item \textbf{Regla 2}: si $p \ge$ techo, fuera (\texttt{0xBF}, vender 63).
\item \textbf{Regla 3}: entre medio, no cambia. Parte fuera.
\item El techo vale 255. Cuando aparece el 255, el techo baja a 225 y se queda ahí.
\item El piso vale 100. Cuando ya aparecieron el 100 y el 255, el piso sube a 120 y se queda ahí.
\item \textbf{Ancla} (solo antes del primer extremo): si $p \ge p_0 + 30$, compra 15 acciones una vez y
  después manda ``mantener''; si $p \le p_0$, vende todo. Toma las subidas del comienzo.
\end{enumerate}
Por qué funciona:
\begin{itemize}
\item Si ve 100 y el 255 todavía no apareció, el 255 viene después: compra en el piso y vende en el
  techo, $\times 2{,}55$ menos comisiones.
\item Si el 255 ya pasó, comprar en 100 no puede perder (el precio no baja de 100). Pero ya no hay
  garantía de volver a 255: por eso toma la ganancia en 225.
\item Cuando ya pasaron los dos, no hay garantía de nada; pero comprar en 120 o menos arriesga poco
  (el precio no baja de 100) y la subida queda abierta.
\end{itemize}
\begin{checkpoint}
Precios: 180, 255, 140, 100, 160, 228, 130, 118. ¿Qué hace? (Fuera; ve 255: el techo baja a 225; fuera;
ve 100: dentro, y como ya vio los dos, el piso sube a 120; sigue dentro en 160; en 228 $\ge 225$ vende;
fuera en 130; en 118 $\le 120$ compra.)
\end{checkpoint}

\subsection*{Paso 4. Ganancia y riesgo}
\addcontentsline{toc}{subsection}{Paso 4. Ganancia y riesgo}
\begin{enumerate}
\item Las rondas son de 20.000 ticks (confirmado). En \NVentanas{} ventanas: media geométrica \ExtGeo{} (CyM \CymGeo), mediana
  \ExtMediana, 5\,\% peor \ExtPcinco, peor ventana \ExtPeor.
\item Pierde plata en el \ExtPierde\,\% de las ventanas. Comprar y mantener: \CymPierde\,\%. El algoritmo
  anterior: \TendPierde\,\%.
\item Parejo: por año 2022 \ExtGeoDosmilveintidos, 2023 \ExtGeoDosmilveintitres, 2024
  \ExtGeoDosmilveinticuatro; mejora en las 5 empresas.
\item Ronda de 2000 ticks: casi siempre el 100 no aparece y no opera. Geo \ExtCortaGeo{} (CyM
  \CymCortaGeo), nunca pierde.
\item Práctica completa: \$\MetaW{} (CyM \$\MetaCym, algoritmo anterior \$\MetaWTend).
\item 100 y 255 salen de la especificación. Techo y recompra se ajustaron (+1,1\,\% y +1,4\,\%); el techo
  quedó en 225 y no en 220 porque le gana cara a cara a reglas parecidas, y cada uno
  mejora a las 5 empresas y a los dos períodos por separado.
\end{enumerate}
\begin{checkpoint}
¿Cuál es el riesgo que queda? (Que el 255 pase primero, compre en 100 y la ventana termine cerca de 100:
peor caso \ExtPeor. Y que la competencia use otra normalización o rondas cortas: ahí no opera y queda en
\$10.000.)
\end{checkpoint}

\subsection*{Paso 5. trade.s, vuelta por vuelta}
\addcontentsline{toc}{subsection}{Paso 5. trade.s, vuelta por vuelta}
\begin{enumerate}
\item Arranque: \texttt{t1} = \texttt{0x8000\_0000}; piso (\texttt{a2}) y mínimo (\texttt{gp}) = \texttt{t1}+100;
  techo (\texttt{a4}) y máximo (\texttt{a3}) = \texttt{t1}+255; \texttt{a0} = \texttt{0x7F}, \texttt{ra} =
  \texttt{0xBF}; \texttt{a5} = \texttt{0xBF} (fuera).
\item \texttt{espera}: \texttt{lw a1, 12(s0)} lee el registro NUEVO: bit 31 = llegó precio, bits 7:0 = precio.
  La lectura es \textbf{bloqueante}: si no hay precio nuevo, el \texttt{lw} no termina (la unidad de carga
  espera), y termina el ciclo siguiente a que llega.
\item Como \texttt{a1} trae el bit 31 en 1, se compara sin signo contra piso y techo, que también lo traen.
\item Lo común: no toca ninguno, escribe DECISION (la misma) y el display, y vuelve.
\item \texttt{piso}: escribe COMPRAR \emph{primero}; después, si $p \le 100$, marca \texttt{s1} y, si ya vio el
  255, sube el piso a 120.
\item \texttt{techo}: escribe VENDER \emph{primero}; después, si $p \ge 255$, baja el techo a 225 y, si ya vio el
  100, sube el piso a 120.
\end{enumerate}
\begin{checkpoint}
¿Por qué escribir DECISION antes de actualizar el estado? (La maestra toma DECISION al bajar CS. Con ticks
pegados quedan ~25 ciclos; actualizar antes hacía llegar tarde justo cuando la decisión cambia: 34 ciclos.
Ahora el peor caso es 9.)
\end{checkpoint}

\subsection*{Paso 6. Cómo se comprobó}
\addcontentsline{toc}{subsection}{Paso 6. Cómo se comprobó}
\begin{enumerate}
\item \texttt{ref/trader.py} (clase \texttt{Extremos}): el algoritmo en Python. Lee las \texttt{.equ} de
  \texttt{trade.s}.
\item \texttt{ref/espino\_iss.py}: ejecuta el \texttt{.hex} instrucción por instrucción contando ciclos.
  En \ExtNTicks{} ticks, misma decisión.
\item \texttt{sim/tb\_trade.v}: el RTL contra una maestra SPI emulada, los 20.000 ticks de la práctica.
  Mismos bytes y \textbf{mismos ciclos por tick que el simulador}.
\item \texttt{make tb}: testbench del curso, \TbAprobados. \texttt{make tb-meta}: el mismo con la
  práctica nueva completa, \TbMetaAprobados, \$\TbMetaW.
\item En la placa: maestra embebida con 1000 precios (tramo con el 255 y después tramo con el 100). 4 LEDs
  fijos y \texttt{A1} en los displays = todo coincide.
\end{enumerate}

\subsection*{Paso 7. Dónde se van los ciclos y Amdahl (algoritmo anterior)}
\addcontentsline{toc}{subsection}{Paso 7. Dónde se van los ciclos y Amdahl}
El estudio de co-diseño se hizo con el algoritmo anterior (\texttt{sw/trade\_tendencia.s}: medias
móviles de 64 y 512, 3 a 5 multiplicaciones por tick). Se reproduce con \texttt{ALGO=tendencia}.
\begin{enumerate}
\item Perfil solo software: multiplicación \FracMul\,\%, duplicaciones \FracDup\,\%.
\item Amdahl: $S = 1/((1-f) + f/s)$. Máximo si la multiplicación costara 0: \AmdahlMaxMul.
\item Predicción antes de construir \texttt{mul8}: cada llamada (unos \CiclosMulSw{} ciclos) pasa a 10.
  $S = \AmdahlPredMul$. Medido: \AmdahlMedMul. Sale mejor porque los umbrales son números chicos y el
  multiplicador termina en 3 a 5 ciclos.
\item Shifter: $f$ = \AmdahlFShift\,\%, $S$ = \AmdahlMedShift. Poco, porque las duplicaciones pesan poco.
\end{enumerate}

\subsection*{Paso 8. El hardware}
\addcontentsline{toc}{subsection}{Paso 8. El hardware}
\begin{enumerate}
\item \textbf{\texttt{no\_rw\_check}} (\texttt{pochoco\_ram.v}): Yosys agregaba 117 flip-flops y 93 LUTs
  para emular leer y escribir la misma dirección en el mismo ciclo. Eso nunca pasa aquí. El SoC bajó de
  \LCOriginal{} a \LCSw{} celdas. Sin esto no cabía ningún acelerador.
\item \textbf{Shifter} (\texttt{espino\_alu.v}, \texttt{ESPINO\_SHIFTER}): un barrel shifter a la derecha
  de 33 bits; SLL invierte los bits antes y después. +\CostoShift{} celdas, $f_\text{máx}$ baja a
  \FmaxShift\,MHz, speedup \SpeedupShift.
\item \textbf{mul8} (\texttt{espino\_mul8.v}): \texttt{rd = rs1 * rs2[7:0]}, opcode custom-0
  \texttt{0x0B}. Secuencial, un bit por ciclo, un solo sumador. Detiene el pipeline como un load.
  +\CostoMul{} celdas, $f_\text{máx}$ igual (\FmaxMul\,MHz), speedup \SpeedupMul. Solo con
  \texttt{-DESPINO\_MUL8}; sin ese define el core es el original.
\item \textbf{Las dos juntas}: \LCMulShift{} celdas, no caben en las 1280.
\item \textbf{PLL}: el VQ100 de la Go Board no tiene (solo el TQ144). nextpnr: ``no suitable BEL''.
\end{enumerate}

\textbf{mul8 ciclo a ciclo} (multiplicar $a$ por $b = 5 = 101_2$):
\begin{center}\small
\begin{tabular}{llll}
\toprule
Ciclo & Qué pasa & acc & stall \\
\midrule
1 & la instrucción llega; espera el banco de registros (BRAM) & -- & sí \\
2 & primer paso: $b_0 = 1$, acc $= a$ & $a$ & sí \\
3 & $b_1 = 0$, no suma; $a \to 2a$ & $a$ & sí \\
4 & $b_2 = 1$, acc $= a + 4a$; no quedan bits en 1: \texttt{done} & $5a$ & sí \\
5 & \texttt{busy} baja: la etapa ID escribe acc en \texttt{rd} & $5a$ & no \\
\bottomrule
\end{tabular}
\end{center}
\begin{checkpoint}
¿Por qué secuencial y no combinacional? (Un $32\times 8$ combinacional son 7 sumadores de 32 bits, unas
250 celdas, y la HX1K no tiene DSP. El secuencial usa 1 sumador.)
\end{checkpoint}

\subsection*{Paso 9. Roofline, Ops/s y la variante final}
\addcontentsline{toc}{subsection}{Paso 9. Roofline, Ops/s y la variante final}
\begin{enumerate}
\item Operación útil = suma, resta, comparación o multiplicación sobre datos del mercado (contadas en
  \texttt{trader.py}): \OpsTick{} por tick, igual en todas las variantes.
\item Techos de la clase: cómputo 12,5\,MOPS (1 op cada 2 ciclos); memoria 100\,MB/s (4 B por ciclo).
  Se cruzan en 0,125 op/byte.
\item El kernel mueve \BytesTick{} bytes por tick: \IntensidadOp{} op/byte. A la derecha del cruce:
  \textbf{limitado por cómputo}.
\item Solo software: \MIPSSw{} MIPS pero \MopsSw{} MOPS útiles. Con \texttt{mul8}: \MopsMul.
\item Algoritmo final: \ExtOpsTick{} ops por tick, \ExtBytesTick{} bytes: 0,25 op/byte, también a la
  derecha del cruce. \ExtMops{} MOPS útiles: casi todas sus instrucciones son de entrada y salida.
\item \textbf{Variante final}: el Espino sin \texttt{mul8} ni shifts, con la lectura bloqueante (985 celdas).
  Es lo que se compila sin opciones.
\item El algoritmo nuevo hizo el tick \ExtSpeedupMul{} veces más corto que el anterior con \texttt{mul8}:
  cambiar el algoritmo aceleró más que el hardware.
\end{enumerate}

%=============================================================================
\section*{Parte B. Preguntas probables}
\addcontentsline{toc}{section}{Parte B. Preguntas probables}

\subsection*{Algoritmo}
\preg{¿Qué predice su algoritmo?}
Nada del mercado: usa cómo se arma la serie. Como el mínimo de la ventana vale 100 y el máximo 255, un
precio de 100 es el piso y un 255 es el techo. Si se ve el piso antes que el techo, el techo viene
después.

\preg{¿No es hacer trampa?}
No. Usa solo los precios que llegan por SPI y la especificación pública de la serie (normalizada entre
100 y 255). Cualquier grupo podía leerla igual.

\preg{¿Por qué no simplemente comprar y mantener?}
Comprar y mantener pierde en el \CymPierde\,\% de las ventanas y su 5\,\% peor deja \CymPcinco{} de la
plata. El agente pierde en el \ExtPierde\,\% y su peor ventana deja \ExtPeor.

\preg{¿Cómo eligieron los parámetros?}
100 y 255 vienen de la especificación. 225 y 120 se eligieron con datos, en \NVentanas{} ventanas, y se
aceptaron solo porque mejoraban a las 5 empresas y a los dos períodos por separado. Para el 120, todo
el rango 115--140 da casi lo mismo: se tomó el medio, no el mejor punto.

\preg{¿Por qué el ancla compra solo 15 acciones?}
Porque antes del primer extremo no hay garantía. Comprando todo, una serie que sube y baja sin tocar
los extremos le hacía perder 20\,\% por vuelta (peor ronda $-21$\,\%). Con 15 acciones la peor ronda
queda en $-5{,}5$\,\% y sigue ganando cara a cara: +8,4\,\% neto contra la regla sin ancla y +12\,\%
contra los rivales.

\preg{¿Por qué 225 y no 220, si 220 da mejor promedio?}
Porque se compite contra grupos con reglas parecidas: entre reglas casi iguales casi todas las rondas
empatan y decide el cara a cara. 225 promedia casi igual (1,997 contra 1,998) y contra un rival con 220
gana en más rondas de las que pierde (+6,2\,\% neto); contra 100/255 puro rinde igual. Se repite en los
ocho escenarios del Montecarlo.

\preg{¿Cómo deciden $q$?}
Siempre 63. La maestra recorta la orden y no cobra si queda en 0. Cerca de 100 alcanza para unas 99
acciones: entra en dos ticks (63 y 36).

\preg{¿Qué pasa si la ronda es de 2000 ticks?}
Casi nunca aparece el 100, así que no opera y termina con \$10.000. Nunca pierde, y en media geométrica
le gana a comprar y mantener (\ExtCortaGeo{} contra \CymCortaGeo).

\preg{¿Qué pasa si no alcanzan a responder?}
La maestra repite la orden anterior, que es la misma que se quería mandar.

\preg{¿Qué pasó con el algoritmo de medias móviles?}
Quedó en \texttt{sw/trade\_tendencia.s}. En las ventanas de la competencia empata con comprar y mantener
(\TendGeo) y pierde en el \TendPierde\,\% de ellas. Con él se hizo el estudio de co-diseño.

\subsection*{Assembly}
\preg{¿Cómo comparan $p \le 100$ sin \texttt{ble}?}
\texttt{bge a2, a1, piso} con \texttt{a2} = 100: es $100 \ge p$.

\preg{¿Usan multiplicación, shifts o división?}
No. Solo comparaciones. Por eso no necesita \texttt{mul8} ni shifts.

\preg{¿Por qué CPI de 2,4 y no 1?}
El banco de registros es una BRAM con lectura sincrónica: cada instrucción espera un ciclo a sus
operandos. Los loads esperan otro ciclo y los saltos tomados descartan la instrucción ya pedida.

\preg{Cambie una constante y regenere el programa.}
Editar la \texttt{.equ} en \texttt{sw/trade.s} (por ejemplo \texttt{TOMA\_GANANCIA}) y:
\begin{lstlisting}
make hex          (sw/trade.hex)
make tb           (el testbench del curso)
make bin && make prog-d2xx
\end{lstlisting}
\texttt{trader.py} lee la misma constante, así que \texttt{python ref/verificar.py --rapido} sigue sirviendo.

\subsection*{Hardware}
\preg{Si el programa final no multiplica, ¿para qué sirvió \texttt{mul8}?}
Para el algoritmo anterior, que sí multiplicaba: bajó su tick de \CiclosPromSw{} a \CiclosPromMul{}
ciclos. Con el final, \texttt{mul8} solo gastaría área, así que el diseño de competencia es el Espino
original. Queda en el repositorio con \texttt{-DESPINO\_MUL8}.

\preg{¿Por qué no un multiplicador más rápido (radix 4)?}
Con el algoritmo final hay cero multiplicaciones. Con el anterior era cerca del 13\,\% de los ciclos:
Amdahl da como mucho 1,07.

\preg{¿Por qué no varios núcleos?}
No caben: el SoC usa \LCSw{} de 1280 celdas y quedan \LCLibres. Recortando el Espino a las
instrucciones que usamos quedaría en unas 400 LUT, pero cada núcleo necesita 4 BRAM para sus registros:
4 núcleos son unas 1600 LUT y 16 BRAM más las 8 de la RAM. Y no hay qué repartir: un precio por tick y
comparaciones que dependen una de otra. Fracción paralela 0, speedup 1, misma plata.

\preg{¿Y con una GPU (CUDA)?}
CUDA corre en GPUs NVIDIA, no en la FPGA. Serviría para buscar parámetros fuera de línea, pero 220
millones de ticks se simulan en 4 segundos en la CPU. Lo que mejoró la ganancia fue entender la
normalización.

\preg{¿Qué hicieron para la prueba de velocidad?}
Tres cosas: escribir DECISION antes de actualizar el estado; un registro NUEVO en el esclavo SPI con
el aviso y el precio en una lectura; y una lectura bloqueante en la unidad de carga del Espino (el
\texttt{lw} de NUEVO no termina hasta que llega un precio). Peor caso precio $\to$ DECISION: de 34 a 9
ciclos (0,36\,$\mu$s), con la misma plata.

\preg{¿Por qué no 4 núcleos o 4 hilos, como otro grupo?}
Llega un precio por tick y la decisión son comparaciones encadenadas: no hay nada que repartir. Un
núcleo con 4 hilos intercalados ejecuta cada hilo a un cuarto de la velocidad, así que la respuesta a
un precio sería más lenta. Sirve para hacer más trabajo independiente, no para responder antes.

\preg{¿Dónde está el cuello de botella ahora?}
En el protocolo, igual para todos: la decisión del precio $n$ se ejecuta al precio $n+1$ y $q \le 63$
obliga a dos órdenes. El cómputo usa el \ExtPctRtres\,\% de la ronda 3.

\subsection*{RTL}
\preg{¿Cómo se integra \texttt{mul8} al pipeline?}
El decoder reconoce el opcode \texttt{0x0B} y marca \texttt{is\_mul} y \texttt{WD\_MUL}. En la etapa ID,
con los operandos listos, arranca \texttt{espino\_mul8}. Su \texttt{busy} se suma a \texttt{lsu\_busy},
así que el controlador detiene el pipeline igual que en un load. Cuando termina, el resultado se
escribe en \texttt{rd}.

\preg{¿Qué es custom-0?}
Un opcode (\texttt{0001011}) que la especificación de RISC-V reserva para extensiones propias.

\preg{¿Sigue siendo Espino?}
Sí. El de la competencia tiene las mismas instrucciones del Espino; solo su unidad de carga acepta una
espera del esclavo SPI. La variante con \texttt{mul8} agrega una instrucción; el decoder, las etapas y
el banco de registros son los originales.

\preg{¿Por qué el shifter baja la frecuencia y \texttt{mul8} no?}
El shifter es combinacional y queda en serie con la ALU. \texttt{mul8} hace un paso por ciclo con un
sumador de 32 bits, tan largo como el de la ALU que ya estaba.

\preg{¿Qué hace \texttt{no\_rw\_check} y por qué es seguro?}
Le dice a Yosys que no emule leer y escribir la misma dirección en el mismo ciclo. El programa no se
escribe a sí mismo y el core no lee y escribe un dato en el mismo ciclo.

\preg{¿Por qué no subieron la frecuencia con el PLL?}
La HX1K en encapsulado VQ100 no tiene PLL. Está en \texttt{resultados/pll\_vq100.txt}.

\subsection*{Métricas}
\preg{¿Dónde está su kernel en el roofline?}
Los dos a la derecha del cruce de 0,125 op/byte: el anterior en \IntensidadOp{} y el final en 0,25.
Limitados por cómputo. Por eso en el anterior se aceleró la multiplicación y no la memoria.

\preg{¿El speedup medido coincide con Amdahl?}
Predicho \AmdahlPredMul, medido \AmdahlMedMul. La predicción supuso el peor caso de \texttt{mul8}.

\preg{¿Cuánto tarda su peor tick?}
\ExtCiclosPeor{} ciclos con el algoritmo final (\ExtPctRtres\,\% de la ronda 3). El anterior:
\CiclosPeorMul{} con \texttt{mul8} y \CiclosPeorSw{} solo software.

%=============================================================================
\section*{Parte C. Números para saberse}
\addcontentsline{toc}{section}{Parte C. Números para saberse}
\begin{center}\small
\begin{tabular}{p{6.2cm}p{9.3cm}}
\toprule
Regla & dentro si $p \le$ piso (100, luego 120); fuera si $p \ge$ techo (255, luego 225) \\
Ventanas evaluadas & \NVentanas{} (META, AAPL, NFLX, AMD, TSLA, 2022--2024) \\
20.000 ticks: geo / mediana / peor & \ExtGeo{} / \ExtMediana{} / \ExtPeor{} (CyM \CymGeo) \\
Prob. de perder (agente, CyM, anterior) & \ExtPierde\,\%, \CymPierde\,\%, \TendPierde\,\% \\
2000 ticks: geo (CyM) & \ExtCortaGeo{} (\CymCortaGeo) \\
Práctica nueva completa: agente / CyM & \$\MetaW{} / \$\MetaCym \\
Testbench: serie por defecto; práctica nueva & \TbAprobados; \TbMetaAprobados, \$\TbMetaW \\
Programa final: palabras; ciclos prom./peor & \ExtPalabras; \ExtCiclosProm{} / \ExtCiclosPeor{} (de \Presupuesto) \\
Celdas del diseño final; $f_\text{máx}$ & \LCSw{} de 1280; \FmaxSw\,MHz \\
\midrule
Anterior: ciclos/tick (sw, shift, mul8) & \CiclosPromSw, \CiclosPromShift, \CiclosPromMul \\
Anterior: fracción en multiplicación (sw) & \FracMul\,\% \\
Amdahl mul8: predicho / medido / máximo & \AmdahlPredMul{} / \AmdahlMedMul{} / \AmdahlMaxMul \\
Celdas: P2, sw, shift, mul8, las dos & \LCOriginal, \LCSw, \LCShift, \LCMul, \LCMulShift{} (de 1280) \\
\bottomrule
\end{tabular}
\end{center}

\end{document}
